% ===================================================================
%  圖表第 11 號 — 城市同佢啲建築。--- 火警。
%
%  Traduction, faite en 2026, en cantonais ecrit d'aujourd'hui, en
%  caracteres traditionnels, sur l'ido de texto/io. Regles de la
%  colonne : voir 10-toubiu-01.tex. Deux parties, deux numerotations :
%  les cles portent c1 et c2.
%
%  LA VILLE EST L'ENDROIT OU HONG KONG A SES PROPRES INSTITUTIONS, ET
%  DONC SES PROPRES MOTS. Le tableau en aligne plus qu'aucun autre :
%
%    差人 (64)     l'agent      Pekin : 警察
%    警司 (94)     le commissaire — le grade tel qu'il s'ecrit
%    郵差 (44)     le facteur   Pekin : 邮递员
%    看更 (37)     le gardien   Pekin : 看守
%    報紙檔 (55)   le kiosque   Pekin : 报刊亭
%    火警 (74)     l'incendie   Pekin : 火灾
%    戲院 (73)     le theatre   Pekin : 剧院
%    收數佬 (33)   l'encaisseur Pekin : 收款员
%    水喉 (69)     le tuyau     Pekin : 水管
%
%  ET LE TROTTOIR (46) DONNE LE CAS LE PLUS NET DE TOUTE LA COLONNE :
%  行人路 contre 人行道. Ce sont les MEMES trois caracteres — homme,
%  marcher, voie — dans un autre ordre, et chaque cote lit le sien
%  sans hesiter. Pas un emprunt, pas une invention : une syntaxe.
%
%  « 差人 » vient de 差 « envoyer en mission » et non de 警 : c'est
%  le mot du XIXe siecle cantonais, celui de 差館 le poste de police,
%  et il a survecu a tous les changements de force publique.
%
%  DEUX BLOCS ONT DEMANDE L'APPOSITION : le dome de l'hopital
%  (c1-12-1), les musiciens de l'orchestre (c2-06-1) ; et deux
%  l'inversion : la porte du chateau (c1-03-1), le foyer du theatre
%  (c2-05-1).
%
%  LE FAC-SIMILE DONNE DEUX FOIS LE RENVOI (81) au bloc c2-06-1 —
%  aux musiciens, puis au violon. La traduction le rend deux fois :
%  la transcription ne corrige pas, et une traduction vise les memes
%  objets dans le meme ordre.
% ===================================================================

%%K t11-apar-1 apar
\VUsaut{2.0mm}
\VUcentre{13.2pt}{90}{第三組}
\VUsaut{3.0mm}
\VUfilet{16mm}
\VUsaut{5.6mm}
\VUtitre{15.0pt}{60}{圖表第 11 號}
\VUsaut{1.4mm}
\VUpk{11.4pt}{40}{城市同佢啲建築。--- 火警。}
\VUsaut{2.2mm}

%%K t11-tit-1 sub
\VUpk{10.2pt}{40}{城邊。}
\VUsaut{3.0mm}

%%K t11-c1-01-1 p
1. --- 我住喺一個大城市度，離大海一百公里，但係佢一樣係個一流嘅\VUgras{港口}
\textsuperscript{(1)}。城邊嗰條大河有四百米闊，仲圍出一個好靚嘅錨地。

%%K t11-c1-02-1 p
2. --- 城入面靠\VUgras{碼頭} \textsuperscript{(2)} 嗰一大片，成個樣都係海同工廠，
變成一個\VUgras{城邊區} \textsuperscript{(3)}，住嘅淨係船員同工人。呢啲工人喺啲
\VUgras{工廠} \textsuperscript{(4)} 同個\VUgras{煤氣廠} \textsuperscript{(5)} 度做嘢；煤氣廠嗰條高
\VUgras{煙囪} \textsuperscript{(6)} 同個\VUgras{氣鼓，}老遠都望得到。

%%K t11-c1-03-1 p
3. --- 中世紀嗰陣，呢個城係有城牆嘅，而家仲搵得返幾段殘跡：尤其係度
\VUgras{城門} \textsuperscript{(8)}——舊\VUgras{城堡} \textsuperscript{(7)} 嗰度——兩邊夾住佢嗰兩座
\VUgras{塔} \textsuperscript{(9)}，而家改咗做\VUgras{監獄。}

%%K t11-c1-04-1 p
4. --- 成個\VUgras{區}望落都好舊，得一座建築算係比較新嘅：就係個\VUgras{海關}
\textsuperscript{(10)}。佢就喺碼頭同條\VUgras{橫街} \textsuperscript{(11)} 之間，條街通去舊
\VUgras{市政廳} \textsuperscript{(12)}。

%%K t11-c1-04-1 p suite
呢座建築好易認：佢有座壓住成個舊城嘅巨型\VUgras{鐘樓} \textsuperscript{(13)}。

%%K t11-c1-05-1 p
5. --- 街嘅另一邊，\VUgras{企}住一座同海關同一個式樣嘅樓：就係\VUgras{交易所}
\textsuperscript{(14)}。佢其中一道門面向住個細廣場，廣場上就係\VUgras{專區公署}
\textsuperscript{(15)}。講真，我哋個城雖然唔係首府，都算係個重要嘅省城。

%%K t11-tit-2 sub
\VUsaut{5.0mm}
\VUpk{10.2pt}{40}{城入面最高嗰一帶。}
\VUsaut{3.4mm}

%%K t11-c1-06-1 p
6. --- 駐軍人數唔少，住喺幾座又闊落又衛生嘅\VUgras{兵營} \textsuperscript{(16)} 度；
營房一路伸到\VUgras{市立公園} \textsuperscript{(17)} 嘅欄杆邊。通去呢個公園嘅條
\VUgras{大馬路} \textsuperscript{(19)} 經過\VUgras{儲蓄銀行} \textsuperscript{(18)} 後面，一直去到
\VUgras{大教堂} \textsuperscript{(20)} 個廣場。

%%K t11-c1-07-1 p
7. --- 呢啲建築好似圓形劇場噉排喺個小山丘上面；我哋間好大嘅\VUgras{中學}
\textsuperscript{(21)} 都喺嗰度。

%%K t11-c1-07-2 p
沿住條斜斜哋、接返大街嘅路行落去，就到\VUgras{法院} \textsuperscript{(22)}；法院前面
攤開一個舒舒服服嘅\VUgras{公園} \textsuperscript{(26)}。呢個\VUgras{花園廣場}唔算大，
但係喺城中心就做咗一片綠色同清涼嘅綠洲。所以市民好鍾意嚟，坐喺\VUgras{茶座}
\textsuperscript{(25)} 嘅帳篷底下，聽軍樂隊逢禮拜四同禮拜日喺個\VUgras{音樂亭}
\textsuperscript{(27)} 度演奏——個亭就擺喺草地同花叢中間。

%%K t11-c1-08-1 p
8. --- 其中一片草地上面企住一座銅\VUgras{像} \textsuperscript{(28)}，係為咗紀念我哋
一位同鄉而立嘅——嗰位為出世嘅城做過好多嘢。

%%K t11-tit-3 sub
\VUsaut{4.0mm}
\VUpk{10.2pt}{40}{交通。}
\VUsaut{3.4mm}

%%K t11-c1-09-1 p
9. --- 到而家為止，我哋個城仲未用電燈照明，淨係得啲\VUgras{煤氣燈}
\textsuperscript{(29)}。不過\VUgras{本地啲報紙}一路喺度催，要改善呢套照明；
\VUgras{就好似}人哋已經改好咗\VUgras{電車嘅牽引方式}噉。啲\VUgras{馬拉}嘅舊
\VUgras{車，}實際上已經畀\VUgras{電車} \textsuperscript{(31)} 取代咗；電車好平就將乘客由
\VUgras{城}嘅一頭\VUgras{送到另一頭。}

%%K t11-c1-10-1 p
10. --- 法院隔籬就係\VUgras{銀行} \textsuperscript{(32)}，商業票據就喺嗰度貼現。啲
\VUgras{收數佬} \textsuperscript{(33)} 就上門去收數。

%%K t11-tit-4 sub
\VUsaut{4.0mm}
\VUpk{10.2pt}{40}{藝術同教育。}
\VUsaut{3.4mm}

%%K t11-c1-11-1 p
11. --- 隔籬企住嗰座靚樓，就係\VUgras{博物館。}入面同時有全城嘅\VUgras{圖書館}
\textsuperscript{(34)}、一批靚到極嘅\VUgras{自然科學藏品} \textsuperscript{(35)}（自然史館），
仲有一間雅雅致致嘅\VUgras{畫館} \textsuperscript{(36)}，本身就係個一流嘅長設\VUgras{展館。}

%%K t11-c1-11-2 p
博物館一個禮拜開兩日畀公眾入；不過外地人畀啲貼士個\VUgras{看更}
\textsuperscript{(37)}，就日日都入得。啲\VUgras{畫家} \textsuperscript{(38)} 會嚟呢度做嘢，
你會見到佢哋\VUgras{企}喺畫架前面，手上攞住\VUgras{調色板，}臨摹啲名畫。

%%K t11-c1-12-1 p
12. --- 博物館後面嗰片高地，隱隱望到個\VUgras{圓頂} \textsuperscript{(40)}——即係
\VUgras{醫院} \textsuperscript{(39)} 嗰個——仲有\VUgras{師範學校} \textsuperscript{(41)} 嘅樓；
我哋啲市立學校嘅老師就係喺嗰度練出嚟。戲院廣場上面有間\VUgras{郵政局}
\textsuperscript{(43)}，設喺一間\VUgras{私人樓} \textsuperscript{(42)} 入面。

%%K t11-tit-5 sub
\VUsaut{5.0mm}
\VUpk{11.4pt}{40}{火警。}
\VUsaut{3.0mm}

%%K t11-tit-6 sub
\VUpk{10.2pt}{40}{「火燭呀！」}
\VUsaut{3.4mm}

%%K t11-c2-01-1 p
1. --- 一個嚇人嘅消息傳遍成個城：戲院啱啱火燭！我趕到個\VUgras{廣場}
\textsuperscript{(96)} 嗰陣，班人已經逼晒喺\VUgras{行人路} \textsuperscript{(46)} 同條
\VUgras{馬路} \textsuperscript{(47)} 度。啲\VUgras{郵差} \textsuperscript{(44)} 同啲\VUgras{派信嘅}
\textsuperscript{(45)} 由郵政局衝咗出嚟。個\VUgras{電車司機} \textsuperscript{(48)} 要停低架車，
讓市政嘅\VUgras{救護車} \textsuperscript{(50)} 過；車上落嚟個\VUgras{外科醫生}
\textsuperscript{(52)} 同個\VUgras{看護} \textsuperscript{(51)}，去救個\VUgras{傷者}
\textsuperscript{(53)}——人哋用\VUgras{擔架} \textsuperscript{(54)} 抬咗佢埋去個\VUgras{報紙檔}
\textsuperscript{(55)} 隔籬。

%%K t11-c2-02-1 p
2. --- 有隊步兵啱啱操完返嚟，就停低幫手維持秩序，同埋啲\VUgras{騎馬嘅市警}
\textsuperscript{(61)} 一齊。\VUgras{啲騎警}啲馬企起前腳，撥人退後就好過班\VUgras{兵}
\textsuperscript{(57)} 或者班\VUgras{憲兵} \textsuperscript{(63)}，將啲\VUgras{睇熱鬧嘅}
\textsuperscript{(56)} 逼開。何況班憲兵都忙到飛起。其中一個喺度同啲\VUgras{差人}
\textsuperscript{(64)} 講，要將另一個\VUgras{傷者} \textsuperscript{(65)} 抬去邊——嗰個係
勇敢嘅消防員，由條梯度跌落嚟。

%%K t11-c2-03-1 p
3. --- \VUgras{個軍官} \textsuperscript{(66)} 指明嚟緊嗰部\VUgras{蒸汽泵}
\textsuperscript{(67)} 要擺喺邊。等緊呢部夠力嘅機器嗰陣，佢先叫人架起部
\VUgras{手泵} \textsuperscript{(68)}，泵條長\VUgras{水喉} \textsuperscript{(69)} 就將水好似洪水噉
射落火度。

%%K t11-c2-03-1 p suite
六個消防員喺度泵，佢哋另一個同伴就攞住個\VUgras{喉嘴} \textsuperscript{(70)}，將條
\VUgras{水柱} \textsuperscript{(71)} 對正火場中心。

%%K t11-c2-04-1 p
4. --- 戲院另一邊，人哋架起咗條大\VUgras{雲梯} \textsuperscript{(72)}。有咗佢，啲
消防員先埋得近啲火舌——啲火喺一團團黑\VUgras{煙} \textsuperscript{(75)} 中間竄出嚟。

%%K t11-tit-7 sub
\VUsaut{5.0mm}
\VUpk{10.2pt}{40}{英勇嘅事。}
\VUsaut{3.4mm}

%%K t11-c2-05-1 p
5. --- 場\VUgras{火警} \textsuperscript{(74)} 由個休息廳燒起——即係\VUgras{戲院}
\textsuperscript{(73)} 嗰個休息廳——當時人哋喺度排\VUgras{歌劇，}\VUgras{首演}本來
聽日先。好彩戲棚入面得幾個\VUgras{觀眾} \textsuperscript{(76)}。呢幾個可憐人避咗上個
\VUgras{露台} \textsuperscript{(77)}，勇敢嘅\VUgras{消防員} \textsuperscript{(86)} 就用梯同粗纜
上去救佢哋。

%%K t11-c2-06-1 p
6. --- 喺\VUgras{柱廊} \textsuperscript{(78)} 底下——嗰度貼住張講戲嘅\VUgras{海報}
\textsuperscript{(79)}——見到啲\VUgras{樂師} \textsuperscript{(81)} 走出嚟，係\VUgras{樂隊}
\textsuperscript{(80)} 嗰班；佢哋孭埋自己啲樂器：\VUgras{小提琴} \textsuperscript{(81)}、
\VUgras{長笛} \textsuperscript{(82)}、\VUgras{大提琴} \textsuperscript{(83)}、\VUgras{長號}
\textsuperscript{(84)}、\VUgras{鼓} \textsuperscript{(85)}，如此類推。大家仲盡力搶救返
一部分\VUgras{傢俬} \textsuperscript{(87)}。

%%K t11-c2-07-1 p
7. --- \VUgras{啲男演員} \textsuperscript{(89)} 同\VUgras{啲女演員} \textsuperscript{(88)}、
\VUgras{男歌手}同\VUgras{女歌手，}着住成套\VUgras{戲服} \textsuperscript{(90)} 就要走。個男高音
喺度同個\VUgras{記者} \textsuperscript{(92)} 講件事點發生。\VUgras{呢位採訪嘅}一開頭就同
啲當官嘅一齊趕到：\VUgras{專員} \textsuperscript{(91)}、\VUgras{市長} \textsuperscript{(93)}、
\VUgras{警司} \textsuperscript{(94)} 同\VUgras{司法官} \textsuperscript{(95)}——佢哋會查清楚，判定
邊個孭鑊。

\VUsaut{6.0mm}
\VUfilet{22mm}
